\begin{pastexam}[label={px:ncku-109-ee-ds-04}]{成大電機所 109 學年度 第 4 題}
（問答；20 分）

\begin{enumerate}[label=(\alph*)]
\item Give conditions under which quick sort algorithm has the worst case behavior in terms of time
complexity. Explain your answer. (10 points)
\item Which sorting algorithm(s) is (are) stable, bubble sort, quick sort, merge sort, heap sort?
\quad (10 points)
\end{enumerate}
\end{pastexam}

\begin{pastsoln}
\textbf{(a)} 快速排序的時間由每次分割切得多平均決定。 \textbf{最壞情況發生在每次分割選到的基準都是
這段子陣列的最大值或最小值}：分割後一邊有 $m - 1$ 個、另一邊是空的，每一層只少一個元素。 分割
長度 $m$ 的子陣列要 $\Theta(m)$ 次比較，所以
\[
T(n) = T(n - 1) + \Theta(n) = \Theta(n) + \Theta(n - 1) + \cdots + \Theta(1) = \Theta(n^2),
\]
比較次數是 $(n - 1) + (n - 2) + \cdots + 1 = n(n - 1)/2$（資結書第~14.5~節）。 遞迴深度也是
$n - 1$，堆疊空間 $\Theta(n)$。

會讓這件事發生的具體條件：
\begin{itemize}
\item 固定取第一個（或最後一個）元素當基準，而輸入\textbf{已經由小到大排好}或\textbf{由大到小排好}。
例如取第一個、輸入已排好時，每次的基準都是最小值。
\item 用 Lomuto 分割（$A[j] \le x$ 的放左邊），而\textbf{所有鍵都相同}：每次基準右邊都是空的，同樣
切成 $m - 1$ 和 $0$。
\item 更一般地說，只要輸入的排列讓每次選到的基準恰好是極值，就會最壞；選基準的規則固定時，總找得到
這種輸入。 隨機選基準（或三數取中）可以讓這種情形極不可能發生，期望時間 $O(n \log n)$。
\end{itemize}

\textbf{(b)} 穩定的是 \textbf{bubble sort} 和 \textbf{merge sort}。 穩定是指鍵相同的資料排序後保持原來的
先後順序（資結書第~14.1~節）。 下面用 $2_a$、$2_b$ 表示鍵相同、原本 $2_a$ 在前的兩筆。
\begin{itemize}
\item Bubble sort：穩定。 只在相鄰兩個\textbf{嚴格大於}（$A[j] > A[j + 1]$）時才交換，相等的兩筆
不會互換，也不會越過彼此（資結書虛擬碼~14.1）。
\item Quick sort：不穩定。 分割時的交換會把元素長距離搬動。 例：$\langle 2_a, 2_b, 1\rangle$ 用
Lomuto 分割（基準是最後的 $1$），最後把基準和 $A[1] = 2_a$ 交換，得 $\langle 1, 2_b, 2_a\rangle$。
\item Merge sort：穩定。 合併時兩邊相等先取左半邊的（比較寫成 $L[i] \le R[j]$），左半邊的資料原本
就在前面（資結書第~14.4~節）。
\item Heap sort：不穩定。 例：$\langle 2_a, 2_b, 1\rangle$ 已是最大堆積，第一次把根 $2_a$ 換到最後，
剩下 $\langle 1, 2_b\rangle$ 調整成 $\langle 2_b, 1\rangle$，再把 $2_b$ 換到第 2 格，結果
$\langle 1, 2_b, 2_a\rangle$，順序顛倒。
\end{itemize}
\end{pastsoln}
